% TOP_PROBLEM: {"id": 2733, "title": "Kirby Problem 1.74", "category_id": 7, "category": {"id": 7, "name": "topology", "display_name": "Topology", "description": "Properties preserved under continuous deformations.", "slug": "topology", "order_index": 7, "created_at": "2026-07-31T15:26:25.671Z"}, "status": "open", "problem_number": "KP-1.74", "set_id": 11, "statusQualification": "The literal strong inequality is refuted by an unknot summand; its nontrivial version explicitly excludes split unknot summing components. The weaker universal positive-saving question retains all nonempty links and necessarily has c at most 2pi."}
% FIRST_POSTED: 2026-10-01
% ENTRY_KIND: statement_only
% UnsolvedMath ID 2733; source catalogue code KP-1.74
% !TeX program = lualatex
% Definition and sources only; this file is not an LLM solution attempt.
\documentclass[11pt,a4paper]{article}
\usepackage[margin=25mm]{geometry}
\usepackage{fontspec}
\setmainfont{lmroman10-regular.otf}[BoldFont=lmroman10-bold.otf,
 ItalicFont=lmroman10-italic.otf,BoldItalicFont=lmroman10-bolditalic.otf]
\usepackage{amsmath,amssymb,mathtools}
\usepackage{unicode-math}
\setmathfont{latinmodern-math.otf}
\AtBeginDocument{\let\setminus\smallsetminus}
\usepackage{xurl,hyperref}
\hypersetup{colorlinks=true,urlcolor=blue}
\setcounter{secnumdepth}{0}
\setlength{\parindent}{0pt}
\setlength{\parskip}{5pt}
\setlength{\emergencystretch}{3em}
\begin{document}
\section*{Kirby Problem 1.74}\label{2733}
{\small UnsolvedMath ID 2733 · KP-1.74 · Topology · Kirby's Problems in Low-Dimensional Topology}
\subsection{Definitions and mathematical statement}
For an embedded link $L\subset\mathbb R^3$, its thickness $\tau(L)$ is the supremal radius of an embedded normal tube about the whole link. Use $C^{1,1}$ curves of positive thickness and count the total length of all components. Define
\[
\operatorname{Rop}(K)=\inf_{L\in K}\frac{\operatorname{length}(L)}{\tau(L)}
\]
for a link type $K$. Thickness incorporates both local curvature and separation between distinct portions or components; the quotient is invariant under rescaling.

Specify a component in each of two links and form their connected sum by deleting small arcs and joining their endpoints compatibly, without extra knotting. Retain all other components. Consider the proposed bound
\[
\operatorname{Rop}(K_1\#K_2)\le
\operatorname{Rop}(K_1)+\operatorname{Rop}(K_2)-(4\pi-4).
\]
The source's literal quantifier over all links is false: the round unknot $U$ has $\operatorname{Rop}(U)=2\pi$, while $K\#U=K$, and $4\pi-4>2\pi$.
One explicit nontrivial link version restricts both selected summing components to those that are not split unknots. Here a split unknot means a component bounding an embedded disk disjoint from the remaining link. This restriction still permits summing along an unknotted component of a Hopf link. Does the displayed bound hold under this explicit restriction?

The source's weaker question can retain all nonempty knot and link types and all component choices: does some absolute $c>0$ satisfy the same inequality with $4\pi-4$ replaced by $c$ universally? Necessarily $c\leq2\pi$, by the unknot example. The constant is independent of crossing numbers and component counts. These are comparisons of optimal ropelengths, not just lengths arising from one prescribed splicing construction; zero saving does not answer the weaker question.
\subsection{Short English statement}
For nontrivial connected sums, does optimal ropelength save the proposed amount? Is some positive saving uniform even when unknot summands are allowed?
\subsection{Sources}
\begin{itemize}
\item[S1] ULAM.AI and the UnsolvedMath Contributors, supplied problems.json, record 2733 (KP-1.74), Kirby's Problems in Low-Dimensional Topology. Catalogue identity and source transcription; CC BY 4.0.
\url{https://huggingface.co/datasets/ulamai/UnsolvedMath}.
\item[S2] K3: A New Problem List in Low-Dimensional Topology, Problem 1.74. Expanded formulation of the supplied catalogue statement.
\url{https://math.berkeley.edu/sites/default/files/surv-295-ruberman-watermarked-author-pdf.pdf}.
\end{itemize}
\end{document}
